\documentclass[11pt]{article}
\usepackage[margin=1.1in]{geometry}
\usepackage{amsmath,amssymb,amsthm}
\usepackage[hidelinks]{hyperref}
\usepackage{microtype}

\newtheorem{lemma}{Lemma}
\newtheorem{theorem}{Theorem}
\newtheorem{corollary}{Corollary}

\newcommand{\Tr}{\operatorname{Tr}}
\newcommand{\ad}{\operatorname{ad}}
\newcommand{\cl}{\mathrm{cl}}
\newcommand{\op}{\mathrm{op}}

\title{Conserved singlet channels in double-scaled SYK:\\
the bilinear tower cannot propagate at $N=\infty$}

% TODO(author): confirm author line and affiliation before circulating.
\author{Henry Donahue}
\date{August 2026 --- draft, not for distribution}

\begin{document}
\maketitle

\begin{abstract}
In the proposed duality between double-scaled SYK at infinite temperature and
the 't~Hooft model [arXiv:2607.05678], the meson tower is to be matched
against correlators of the singlet bilinears
$A_n=\sum_i\psi_i(\ad_H)^n\psi_i$. We prove that at strict double scaling
($N=\infty$, any $\lambda$) these channels cannot propagate: in the chord
representation, the insertion of $A_n$ into a pair correlator acts as a
\emph{polynomial in the chord Hamiltonian}, so every spectral moment
$\mu_k$, $k\ge1$, of every $A_n$ pair channel vanishes identically --- at
every coupling and every matter weight. The mechanism is a closed three-term
recursion for the ``dressed propagator'' $C(m)=\cl\circ T_M^m\circ\op$
(open a matter chord, apply the one-particle transfer matrix $m$ times,
close it), whose coefficients are scalars; $C(m)$ is the Motzkin-path
polynomial of $T$ with flat weight $q^{\Delta+j}T$, down weight $[j]_q$ and
up weight $1-q^{2\Delta+j}$ at height $j$. The first conserved charges are
$\widehat A_1=-(1-q^{\Delta})\,T$ (the chord image of the exact identity
$A_1=-2pH$) and $\widehat A_2=(1-q^{\Delta})^2T^2+(1-q^{2\Delta})$. By
contrast a random matter operator inserts $q^{\Delta\hat n}$, which does not
commute with $T$ and therefore propagates. Consequently the meson tower of
the proposed duality is a $1/N$ effect, invisible at every order of the
double-scaled expansion; the quantitative spectrum match must be performed
at fixed $p$, $\lambda\to0$, where the emergent gauge symmetry is only
polynomially accurate. Every statement is machine-verified in an
accompanying repository.
\end{abstract}

\section{Setting and statement}

Double-scaled SYK ($N$ Majorana fermions, $H=i^{p/2}\sum J\,\psi^p$,
$\lambda_{\rm chord}=2p^2/N$ fixed, $q=e^{-\lambda_{\rm chord}}$) is solved
at $N=\infty$ by chord combinatorics \cite{BINT,BM}. Ensemble-averaged
traces reduce to amplitudes on the chord Hilbert space
$\{\,|n\rangle\,\}_{n\ge0}$, with the chord Hamiltonian acting (in the
formal basis, units $\langle H^2\rangle=1$) as
\begin{equation}
T\,|n\rangle=|n+1\rangle+[n]_q\,|n-1\rangle,
\qquad [n]_q=\frac{1-q^n}{1-q}.
\end{equation}
A matter operator of dimension $\Delta$ opens a matter chord; while it is
open the Hamiltonian acts by the one-particle transfer matrix $T_M$ on
states $|n_0,n_1\rangle$ ($n_0$ H-chords opened before the matter chord,
$n_1$ after):
\begin{equation}
T_M\,|n_0,n_1\rangle=|n_0,n_1{+}1\rangle+[n_1]\,|n_0,n_1{-}1\rangle
+q^{\Delta+n_1}[n_0]\,|n_0{-}1,n_1\rangle ,
\label{eq:TM}
\end{equation}
the three terms being: open a new H-chord; close one opened after the
matter chord; close one opened before it (crossing the matter chord, weight
$q^{\Delta}$, and the $n_1$ later chords, weight $q^{n_1}$)
\cite{BINT,BM}. The matter chord itself opens and closes via
\begin{equation}
\op\,|n\rangle=|n,0\rangle,
\qquad
\cl\,|n_0,n_1\rangle=q^{\Delta n_1}\,|n_0+n_1\rangle .
\end{equation}

For flavors $i\neq j$, the ensemble average of a \emph{sequential} pair
word factorizes through the zero-chord sector,
\begin{equation}
\Big\langle\Tr\big[\,H^{a}\psi_i H^{b}\psi_i\,H^{c}\psi_j H^{d}\psi_j\big]
\Big\rangle_J
=\langle0|\;T^{a}\,C(b)\;T^{c}\,C(d)\;|0\rangle,
\qquad
C(m):=\cl\circ T_M^{\,m}\circ\op ,
\label{eq:factor}
\end{equation}
with $\Delta=\Delta_\psi$ the single-fermion weight ($\Delta_\psi=1/p$ for
the leading chord weight; nothing below depends on its value). We call
$C(m)$ the \emph{dressed propagator}: it is the two-point insertion
$\psi\,H^m\,\psi$ seen from the zero-chord sector.

The channels of interest are the singlet bilinears of \cite{MSS2607},
\begin{equation}
A_n=\sum_i\psi_i(\ad_H)^n\psi_i,
\qquad
(\ad_H)^n\psi=\sum_{r=0}^{n}(-1)^r\binom{n}{r}H^{\,n-r}\psi H^{\,r}.
\end{equation}
By \eqref{eq:factor} and linearity, inserting $A_n$ (per flavor pair) acts
on the chord line as
\begin{equation}
\widehat A_n=\sum_{r=0}^{n}(-1)^r\binom{n}{r}\,C(n-r)\,T^{\,r}
\;=\;\frac{d^n}{dt^n}\Big[\,C(t)\,e^{-tT}\Big]_{t=0},
\qquad C(t):=\cl\,e^{tT_M}\op ,
\label{eq:Ahat}
\end{equation}
and the spectral moments of the pair channel are
\begin{equation}
\mu_k \;=\; \big\langle\Tr\big[\ad_H^k(X)\,Y\big]\big\rangle_J
\;=\;\langle0|\,\ad_T^k\big(\widehat A_n\big)\,\widehat B_n\,|0\rangle ,
\label{eq:moments}
\end{equation}
where $X=\psi_i(\ad^n\psi_i)$, $Y=X^\dagger$ on flavor $j$, and
$\widehat B_n$ is the chord image of the $Y$ block (its form is irrelevant
below). Our result:

\begin{theorem}[conserved channels]
$C(m)$ is a polynomial in $T$ for every $m$. Hence $\widehat A_n$ is a
polynomial in $T$, $\ad_T^k(\widehat A_n)=0$, and every moment
\eqref{eq:moments} vanishes for every $n\ge1$ and every $k\ge1$ --- at
every $q\in[0,1)$ and every matter weight $\Delta$.
\end{theorem}

\section{The recursion}

For $j\ge0$ define the $j$-fold closes
\begin{equation}
\cl_j\,|n_0,n_1\rangle
=[n_1][n_1{-}1]\cdots[n_1{-}j{+}1]\;q^{\Delta(n_1-j)}\,|n_0+n_1-j\rangle ,
\qquad \cl_0=\cl ,
\end{equation}
(zero if $n_1<j$).

\begin{lemma}[three-term recursion, scalar coefficients]
\begin{equation}
\cl_j\circ T_M
=q^{\Delta+j}\;T\circ \cl_j
\;+\;[j]_q\;\cl_{j-1}
\;+\;\big(1-q^{2\Delta+j}\big)\;\cl_{j+1}.
\label{eq:rec}
\end{equation}
\end{lemma}

\begin{proof}
Apply both sides to $|n_0,n_1\rangle$ and match the two chord-number
sectors. Write $[n_1]^{(j)}=[n_1]\cdots[n_1{-}j{+}1]$ and
$G=n_0+n_1$. Raising sector ($|G-j+1\rangle$): the left side receives only
the open term of \eqref{eq:TM}, giving $[n_1{+}1]^{(j)}q^{\Delta(n_1+1-j)}$;
the right side gives $q^{\Delta+j}[n_1]^{(j)}q^{\Delta(n_1-j)}$ from
$T\,\cl_j$ plus $c_j\,[n_1]^{(j-1)}q^{\Delta(n_1-j+1)}$ from $\cl_{j-1}$.
Using $[n_1{+}1]^{(j)}=[n_1{+}1]\,[n_1]^{(j-1)}$ and
$[n_1]^{(j)}=[n_1]^{(j-1)}[n_1{-}j{+}1]$, the matching condition is
\begin{equation}
c_j=[n_1{+}1]-q^{\,j}\,[n_1{-}j{+}1]
=\frac{(1-q^{n_1+1})-q^j(1-q^{n_1-j+1})}{1-q}=[j]_q ,
\end{equation}
independent of $n_1$. Lowering sector ($|G-j-1\rangle$): the left side
receives the two close terms of \eqref{eq:TM},
$[n_1]^{(j+1)}q^{\Delta(n_1-1-j)}+q^{\Delta+n_1}[n_0][n_1]^{(j)}q^{\Delta(n_1-j)}$;
on the right, $T\,\cl_j$ contributes
$q^{\Delta+j}[n_1]^{(j)}\big([n_1{-}j]+q^{\,n_1-j}[n_0]\big)q^{\Delta(n_1-j)}$
(using $[G-j]=[n_1{-}j]+q^{\,n_1-j}[n_0]$) and $\cl_{j+1}$ contributes
$d_j\,[n_1]^{(j+1)}q^{\Delta(n_1-j-1)}$. The $[n_0]$ terms cancel
identically, and with $[n_1]^{(j+1)}=[n_1]^{(j)}[n_1{-}j]$ the remaining
condition is $q^{-\Delta}=q^{\Delta+j}+d_j\,q^{-\Delta}$, i.e.
$d_j=1-q^{2\Delta+j}$, again independent of $n_1$ (and of $n_0$).
\end{proof}

Iterating \eqref{eq:rec} from $\cl_0$, every application of $T_M$ replaces
$\cl_j$ by a combination of $\cl_{j'}$ with prefactors that are scalars or
scalar multiples of $T$ acting on the left. Since
$\cl_{j}\circ\op=\delta_{j0}$, projecting onto $j=0$ after $m$ steps gives
\begin{equation}
\boxed{\;C(m)=P_m(T)\;}
\end{equation}
with $P_m$ the \emph{Motzkin-path polynomial}: sum over paths
$0\to0$ of length $m$ in the auxiliary height $j$, a flat step at height
$j$ contributing $q^{\Delta+j}\,T$, a down step $[j]_q$, an up step
$1-q^{2\Delta+j}$. Equivalently, $\sum_m P_m(T)\,t^m$ is the $(0,0)$
resolvent of the auxiliary tridiagonal operator --- the continued fraction
\begin{equation}
\sum_{m\ge0}P_m(T)\,t^m=
\cfrac{1}{1-t\,q^{\Delta}T-
 \cfrac{t^2\,(1-q^{2\Delta})[1]}{1-t\,q^{\Delta+1}T-
  \cfrac{t^2\,(1-q^{2\Delta+1})[2]}{1-\cdots}}} \; .
\end{equation}
The first cases are
\begin{equation}
C(0)=1,\qquad
C(1)=q^{\Delta}\,T,\qquad
C(2)=q^{2\Delta}\,T^2+\big(1-q^{2\Delta}\big),
\end{equation}
and from \eqref{eq:Ahat}
\begin{equation}
\widehat A_1=-(1-q^{\Delta})\,T,
\qquad
\widehat A_2=(1-q^{\Delta})^2\,T^2+\big(1-q^{2\Delta}\big).
\end{equation}
$\widehat A_1\propto T$ is the $N=\infty$ image of the finite-$N$ operator
identity $A_1=-2pH$ (exact per disorder instance); $\widehat A_2$ is, to
our knowledge, the first explicit higher conserved charge of the tower.

\section{Corollaries and consequences}

\begin{corollary}[what propagates and what cannot]
A random matter operator $O_\Delta$ (independent Gaussian couplings)
inserts $q^{\Delta\hat n}$, which does not commute with $T$; its channel
propagates, with $\mu_2=2(1-q^{\Delta})$. The $H$-derived bilinears insert
polynomials in $T$ and cannot propagate. The distinction is not the size of
the operator but the correlation of its couplings with those of $H$.
\end{corollary}

\begin{corollary}[the meson tower is a $1/N$ effect]
In the proposed DSSYK$_\infty$--'t~Hooft correspondence
\cite{MSS2607,MSS2506}, the singlet tower $M_n$ maps to the $A_n$ channels.
At strict double scaling their connected pair correlators are
time-independent to all orders of the $\lambda$-expansion: any meson tower
must arise entirely from $1/N$ corrections to the chord solution. The
quantitative spectrum match the authors propose as ``the best test'' must
therefore be performed at fixed $p$ (fixed $\bar\alpha=p^2$), $\lambda\to0$
--- the regime where, as shown in the companion note on exact disorder
statistics, the emergent gauge symmetry is only polynomially accurate
(symmetry violations $\sim N^{-(p-1)/2}$). This sharpens, exactly in
$\lambda$, the large-$q$ observation that the bilinear-tower couplings
first appear at $O(1/q^2)$ \cite{GR}.
\end{corollary}

Two scope remarks. First, the theorem concerns the $i\neq j$ pair
amplitude at strict $N=\infty$ --- precisely the object the chord expansion
computes; finite-$N$ corrections and the flavor-diagonal part of the
physical channel are untouched (numerically, at exact-diagonalization
sizes, the raw channel width grows toward the fixed-$p$ corner). Second,
conservation here is a statement about moments of the connected pair
correlator, $\mu_k=0$ for all $k\ge1$; the weight $\mu_0$ itself is
nonzero.

\paragraph{Verification.}
Every step is machine-verified in the accompanying
repository:\footnote{\texttt{duality/chord\_charges.py} and
\texttt{test\_chord\_charges.py}; the original numerical discovery and the
$n\le5$, $\mu_{2,4,6}$ scans are \texttt{duality/chord\_moments.py}.} the
recursion \eqref{eq:rec} as a matrix identity for $j\le5$; $C(m)=P_m(T)$
and $[C(m),T]=0$ for $m\le8$; the closed forms above; and, through an
independent transfer-matrix evaluation of the trace words (itself validated
against brute-force chord-diagram enumeration), the vanishing of
$|\mu_k|$ relative to its pre-cancellation gross magnitude at the level of
$10^{-17}$ for $n\le5$, $k\le6$, across $\lambda\in[0.1,4]$ and matter
weights $\Delta_\psi\in[0.1,1]$.

\paragraph{Acknowledgments.}
% TODO(author): acknowledgments as appropriate.
This note grew out of a computational stress-test of \cite{MSS2607}; all
code and machine-verified derivations are available from the author.

\begin{thebibliography}{9}

\bibitem{MSS2607}
S.~Miyashita, Y.~Sekino and L.~Susskind,
``Holograms and Standard Models,''
arXiv:2607.05678 [hep-th].

\bibitem{MSS2506}
S.~Miyashita, Y.~Sekino and L.~Susskind,
``DSSYK at infinite temperature: the flat-space limit and the 't~Hooft model,''
JHEP \textbf{11} (2025) 107, arXiv:2506.18054 [hep-th].

\bibitem{BINT}
M.~Berkooz, M.~Isachenkov, V.~Narovlansky and G.~Torrents,
``Towards a full solution of the large $N$ double-scaled SYK model,''
JHEP \textbf{03} (2019) 079, arXiv:1811.02584 [hep-th].

\bibitem{BM}
M.~Berkooz and O.~Mamroud,
``A Cordial Introduction to Double Scaled SYK,''
Rept.\ Prog.\ Phys.\ \textbf{88} (2025) 036001, arXiv:2407.09396 [hep-th].

\bibitem{GR}
D.~J.~Gross and V.~Rosenhaus,
``The bulk dual of SYK: cubic couplings,''
JHEP \textbf{05} (2017) 092, arXiv:1702.08016 [hep-th].

\end{thebibliography}

\end{document}
